% =============================================================================
% Adaptive AI for Patient Reactivation and Causal Impact Measurement
% in Outpatient Clinics -- PFE / Graduation Report
%
% This file was generated from ACADEMIC_REPORT_DRAFT.md and could not be
% compiled locally in the environment that produced it (no TeX toolchain was
% available). It uses only standard, widely-available packages and avoids
% any institution-specific class file, so it should compile as-is on
% Overleaf or any reasonably current TeX Live installation. Please do a
% test compile (pdflatex, twice, then once more for the ToC/labels to
% settle) before relying on the PDF, and fix anything your own template
% requires (a required school-specific .cls file, exact title-page layout,
% logo images, etc.).
%
% Placeholders you must fill in before submission are marked "TODO:".
% =============================================================================

\documentclass[12pt,a4paper]{report}

% ---- Encoding / fonts ----
\usepackage[utf8]{inputenc}
\usepackage[T1]{fontenc}
\usepackage{lmodern}

% ---- Underscores in plain text without manual escaping ----
\usepackage[strings]{underscore}

% ---- Page geometry ----
\usepackage[a4paper,left=3cm,right=2.5cm,top=2.5cm,bottom=2.5cm,headsep=1cm]{geometry}

% ---- Headers / footers ----
\usepackage{fancyhdr}
\pagestyle{fancy}
\fancyhf{}
\renewcommand{\headrulewidth}{0.4pt}
\fancyhead[L]{\small\itshape\leftmark}
\fancyfoot[C]{\small\thepage}
\fancypagestyle{plain}{%
  \fancyhf{}
  \fancyfoot[C]{\small\thepage}
  \renewcommand{\headrulewidth}{0pt}
}

% ---- Chapter / section styling ----
\usepackage{titlesec}
\titleformat{\chapter}[hang]
  {\normalfont\huge\bfseries}
  {\thechapter}
  {20pt}
  {\huge}
\titlespacing*{\chapter}{0pt}{0pt}{30pt}

% ---- Tables, figures, misc ----
\usepackage{booktabs}
\usepackage{array}
\usepackage{longtable}
\usepackage{graphicx}
\usepackage{caption}
\usepackage{amsmath}
\usepackage{amssymb}
\usepackage{enumitem}
\usepackage{xcolor}
\usepackage{tocloft}
\usepackage[colorlinks=true,linkcolor=black,citecolor=black,urlcolor=blue]{hyperref}

% ---- Table of contents depth ----
\setcounter{tocdepth}{3}
\setcounter{secnumdepth}{3}

% ---- Convenience macros ----
\newcommand{\code}[1]{\texttt{#1}}
\newcommand{\university}{TODO: Your University / School}
\newcommand{\program}{TODO: Your Engineering Program / Option}
\newcommand{\authorname}{TODO: Your Name}
\newcommand{\defensedate}{TODO: Defense Date}

\begin{document}

% =============================================================================
% TITLE PAGE
% =============================================================================
\begin{titlepage}
\centering
\vspace*{1cm}

{\small TODO: Ministry / country line \par}
{\small \university \par}
\vspace{0.5cm}
\rule{0.6\textwidth}{0.4pt}
\vspace{1cm}

{\LARGE\bfseries Graduation Report\par}
\vspace{0.5cm}
{\large Presented at\par}
\vspace{0.3cm}
{\Large \university \par}
\vspace{0.8cm}
{\large In order to obtain the\par}
\vspace{0.2cm}
{\Large\bfseries Engineering Diploma in: \program \par}
\vspace{1.2cm}

{\large By\par}
\vspace{0.3cm}
{\LARGE\bfseries \authorname \par}
\vspace{1cm}
\rule{0.4\textwidth}{0.6pt}
\vspace{0.6cm}

{\Large\bfseries Adaptive AI for Patient Reactivation and\\
Causal Impact Measurement in Outpatient Clinics\par}
\rule{0.4\textwidth}{0.6pt}
\vspace{1.2cm}

{\itshape Presented on \defensedate, before the review panel\par}
\vspace{0.8cm}

\begin{tabular}{@{}ll@{}}
\textbf{Mr./Ms.} & TODO: President \\[4pt]
\textbf{Mr./Ms.} & TODO: Examiner \\[4pt]
\textbf{Mr./Ms.} & TODO: Academic Supervisor \\[4pt]
\textbf{Mr.\ Wael Hilali} & Industrial Supervisor \\[4pt]
\textbf{Mr.\ Bilel Said} & Industrial Supervisor \\
\end{tabular}

\vfill
{\small Velodoc \par}

\end{titlepage}

% =============================================================================
% DEDICATION
% =============================================================================
\chapter*{Dedication}
\addcontentsline{toc}{chapter}{Dedication}

TODO: Personalize this page. The convention used across this cohort's
reports dedicates the work to parents, family, friends, the academic
community at \university, and colleagues at the host company, Velodoc.

\vspace{2cm}
\begin{flushright}
\authorname
\end{flushright}

% =============================================================================
% ACKNOWLEDGEMENT
% =============================================================================
\chapter*{Acknowledgement}
\addcontentsline{toc}{chapter}{Acknowledgement}

TODO: Personalize this page. Thank your academic supervisor, your
industrial supervisors at Velodoc (Mr.\ Wael Hilali, Mr.\ Bilel Said),
the members of the jury, the professors and staff of \university, and
your family and friends, following the same structure used across this
cohort's reports.

\vspace{2cm}
\begin{flushright}
\authorname
\end{flushright}

% =============================================================================
% FRONT MATTER: TOC, LOF, LOT
% =============================================================================
\pagenumbering{roman}

\tableofcontents
\listoffigures
\listoftables

\chapter*{List of abbreviations}
\addcontentsline{toc}{chapter}{List of abbreviations}

\begin{longtable}{@{}p{3cm}p{11cm}@{}}
\toprule
AI & Artificial intelligence \\
API & Application programming interface \\
AUC & Area under the curve \\
Brier & Brier score \\
CatBoost & Categorical boosting \\
CDT & Current Dental Terminology (future-work scope only) \\
DGX Spark & NVIDIA DGX Spark \\
ECE & Expected calibration error \\
EHR & Electronic health record \\
FHIR & Fast Healthcare Interoperability Resources \\
HITL & Human in the loop \\
JSON & JavaScript Object Notation \\
LLM & Large language model \\
ML & Machine learning \\
MLOps & Machine learning operations \\
NATS & NATS JetStream (message bus) \\
OAuth & Open Authorization \\
PR-AUC & Precision-recall area under the curve \\
REST & Representational state transfer \\
ROC & Receiver operating characteristic \\
SaMD & Software as a Medical Device \\
SHAP & Shapley additive explanations \\
SMS & Short Message Service \\
\bottomrule
\end{longtable}

\chapter*{General introduction}
\addcontentsline{toc}{chapter}{General introduction}
\pagenumbering{arabic}

Outpatient clinics accumulate, silently, a large population of patients
who have a legitimate clinical or operational reason to come back --- a
treatment plan left open, an insurance benefit about to expire, a
clinical guideline recall, or simple dormancy after a long gap since the
last visit --- but are never proactively recontacted. Front-desk teams
are structured to react to inbound demand (a patient calling to book),
not to mine their own patient population for who \emph{should} be
contacted next; when outreach does happen, it is usually a blanket
campaign with no personalization and no way to prove it changed
anything.

This PFE internship was carried out with \textbf{Velodoc}, a healthcare
AI startup, as part of the same broader company vision described
elsewhere in this cohort's reports: adaptive AI for clinic operations,
spanning scheduling intelligence, patient reactivation, and
clinical-operational copilots. The work presented here is the
\textbf{patient-reactivation and engagement-measurement} track of that
vision, implemented as \textbf{Velo Engage}, branded to clinic users as
\textbf{VeloDoc}.

The system finds patients with a legitimate reason to be recontacted,
checks whether contacting them is legally and operationally permitted,
decides --- with trained machine-learning models, not a fixed rule ---
whether contacting them is actually worth it, reaches them on WhatsApp,
and measures whether the outreach caused a real incremental booking
rather than simply reaching people who would have returned anyway. It
runs once a day, per clinic, and treats a \textbf{randomized holdout
experiment} as a first-class architectural component, not an
afterthought --- the platform is built to prove its own causal impact,
not just report a correlation.

This report is organized into four chapters. Chapter~\ref{ch:context}
presents the host company, the problem, existing approaches and their
limits, the proposed solution, and the requirement analysis.
Chapter~\ref{ch:architecture} presents the system and AI architecture:
the EHR integration, the daily pipeline, the model choices, and the
MLOps/governance design. Chapter~\ref{ch:implementation} describes the
implementation of each service. Chapter~\ref{ch:evaluation} presents the
evaluation methodology and the results obtained, together with the
project's honest limitations. A general conclusion and future
perspectives close the report.

% =============================================================================
\chapter{Adaptive AI for Patient Reactivation and Engagement Measurement}
\label{ch:context}

\section{Introduction}

This chapter presents the internship context, the operational problem
this project addresses, the study and critique of existing patient-recall
approaches, the proposed solution, the underlying concepts, the
requirement analysis, the use case diagram, and the project's planning
and development approach.

\section{Host company and internship context}

\subsection{Host company presentation}

Velodoc is a healthcare artificial intelligence startup founded in 2024,
operating in the hospitals and health care sector, focused on AI systems
embedded directly in healthcare workflows. It addresses administrative
fragmentation in healthcare --- doctors, billing teams, clinic owners,
and front-desk teams working with disconnected tools, causing repeated
work, manual follow-up, limited operational visibility, and delays.
Velodoc develops healthcare workflow copilots and operational
intelligence tools, with the explicit objective of supporting healthcare
teams rather than replacing clinical or operational judgment.

\subsection{Internship mission}

The mission for this internship was the \textbf{patient-reactivation and
engagement-measurement} part of Velodoc's broader vision: build a system
able to (1)~detect, from a clinic's own patient population, every patient
with a legitimate reason to be recontacted; (2)~gate that list through
consent and contact-frequency rules; (3)~rank survivors with trained
machine-learning models rather than a static priority; (4)~dispatch
outreach over WhatsApp with an automatic SMS fallback; (5)~measure the
causal effect of that outreach through a randomized holdout, not just
report correlational outcomes; and (6)~expose all of this through a
role-gated web console for clinic staff, owners, and administrators.

\subsection{Collaboration tools}

The project was developed through the same collaboration setup used
across the Velodoc team: Notion for specifications, GitHub for source and
version control, Discord for daily communication, and a shared DGX Spark
environment for GPU-bound components (the clinical-note extraction
model, described in Chapter~\ref{ch:implementation}). Development also
relied on a live Epic FHIR sandbox for real, non-synthetic EHR
integration testing --- a deliberate choice over building an EHR
emulator from scratch, discussed in Section~\ref{sec:critique}.

\section{Problem statement}

A clinic's own patient population already contains most of the ``leads''
it needs --- patients aren't lost to a competitor, they are simply never
recontacted. Several distinct situations create a legitimate reason to
reach back out: a patient hasn't visited in a long time relative to
their own normal cadence (dormancy); a patient has an \textbf{open,
incomplete treatment plan}; a patient's insurance benefit is about to
expire, with no clinical or financial reason to let the window close
unused; a clinical guideline indicates a recall is due (age, condition,
medication); and several other operational patterns.

Even when a clinic recognizes this, three problems remain unsolved by a
simple reminder system:

\begin{enumerate}
  \item \textbf{No personalization of \emph{whether} to contact.} A
        blanket campaign to everyone with an open reason ignores that
        some patients are highly unlikely to attend even if booked, some
        would attend regardless of any nudge, and some represent much
        higher expected value if they return than others.
  \item \textbf{No enforcement of \emph{who may legally be contacted
        about what}.} Consent for one type of contact (e.g., a routine
        recall) does not imply consent for another (e.g., a promotional
        offer) --- a single flat ``opted in'' flag is not enough.
  \item \textbf{No proof of causal impact.} A clinic that sees bookings
        rise after a campaign cannot distinguish ``the campaign worked''
        from ``these patients would have booked anyway'' without a
        randomized comparison group.
\end{enumerate}

The central problem this project addresses: \emph{how can a system
detect legitimate re-engagement opportunities across a clinic's own
patient population, gate them through real consent and contact rules,
rank them by trained models of expected value and likely response,
dispatch outreach automatically, and prove --- not assume --- that the
outreach caused incremental bookings?}

\section{Study of existing solutions}

\textbf{Manual chart review and call lists.} The traditional approach:
staff periodically review charts or run ad hoc reports to find patients
due for follow-up, then call them. This does not scale past a small
patient panel and depends entirely on staff bandwidth.

\textbf{Generic recall/reminder systems built into EHR or
practice-management software.} Most EHR platforms support a basic recall
flag (e.g., ``due for a 6-month cleaning'') and a bulk SMS/email blast.
This handles one narrow recall type well but does not unify dormancy,
open treatment plans, benefit expiry, and clinical-guideline recall into
a single ranked list, and does not personalize contact order by
predicted outcome.

\textbf{Mass marketing/CRM tools repurposed for healthcare.}
General-purpose marketing automation platforms can segment and blast
messages, but they are not built around per-contact-type consent, are
not integrated with a clinical EHR as source of truth, and have no
notion of a causal holdout to measure real impact versus noise.

\textbf{Academic and industry work on patient no-show/engagement
prediction} exists (as surveyed extensively in the sister PFE report on
no-show prediction and slot recovery from this same team), but that
literature is concentrated on \emph{predicting attendance for an
already-booked appointment}, not on the earlier, upstream problem this
project addresses: \emph{deciding which never-booked patient to
proactively reach out to in the first place, and proving the outreach
caused the booking.}

\section{Critique of existing solutions}
\label{sec:critique}

Existing recall tooling is narrow (one recall type at a time),
consent-blind (a single opt-in flag rather than per-purpose consent
classes), and unranked (patients are contacted in registration order or
by a fixed static priority, never personalized by a trained model of
expected outcome). Most critically, none of the surveyed approaches
build in a mechanism to \emph{prove} the campaign caused the observed
bookings --- without a randomized control group, a rise in bookings
after a campaign is only ever a correlation.

\section{Proposed solution}

The proposed solution is \textbf{Velo Engage}, branded to clinic end
users as \textbf{VeloDoc}. It runs, once per clinic per day, a
five-stage pipeline: \textbf{detect $\rightarrow$ gate $\rightarrow$ rank
$\rightarrow$ dispatch $\rightarrow$ measure}.

\begin{itemize}
  \item \textbf{Detect}: 16 independent, YAML-defined ``opportunity
        family'' rules evaluate every patient --- dormancy, open
        treatment plans, benefit expiry, clinical recall (backed by a
        knowledge graph), and others --- each carrying its own
        consent-class requirement and priority weight. One family is not
        rule-based at all: it uses a lookalike/collaborative-filtering
        model to suggest a family a patient has never used, based on
        similar patients who have.
  \item \textbf{Gate}: candidates are filtered down to patients who have
        specifically consented to \emph{that} type of contact, have not
        been messaged too recently or too often, and have their
        multiple simultaneous reasons for contact collapsed into one.
  \item \textbf{Rank}: five trained machine-learning models --- no-show
        risk, booking propensity, expected revenue, causal uplift, and
        predicted time-to-book --- combine into a single expected-value
        ranking. A randomly-selected slice of the ranked population is
        deliberately \textbf{held back from contact entirely} (the
        holdout arm), regardless of rank, to measure the program's true
        causal effect.
  \item \textbf{Dispatch}: the highest-ranked, non-holdout patients (a
        deliberately capped number per day) receive a WhatsApp message,
        with automatic SMS fallback if delivery fails.
  \item \textbf{Measure}: replies, bookings, attendance, and revenue flow
        back automatically via webhooks, or are recorded manually when
        no automated feed exists yet, closing the loop for the next
        day's ranking and for eventual model retraining. Recovered
        revenue is weighed against outreach cost to produce a real ROI
        figure.
\end{itemize}

The EHR (Epic, via its FHIR API) remains the system of record for
clinical data; Velo Engage reads from it, computes its own
operational/predictive data in its own Postgres schema, and only ever
writes back real, already-approved outreach actions (a WhatsApp message,
a booking record) --- never a clinical fact.

This implemented scope explicitly does \textbf{not} claim: production
Epic credentials for a real (non-sandbox) clinic, formal Meta WhatsApp
Business template approval at commercial scale, or a fully-reviewed
clinical decision-support extension into treatment-pathway
recommendation (a prototype for this last item is described as future
work in Section~\ref{sec:dental-prototype}).

\section{Basic concepts}

\subsection{Opportunity detection}

The mechanism that decides \emph{who} has a legitimate reason to be
recontacted. Implemented as YAML-defined rules rather than code
specifically so a non-engineer can review and edit the logic --- most
families are deterministic (thresholds on dormancy days, benefit-expiry
windows), one is backed by a Neo4j clinical-guideline knowledge graph,
and one is backed by a trained lookalike model.

\subsection{Consent and guardrails}

Per-class consent (contacting a patient about a clinical recall does not
imply consent for a promotional offer), contact cooldowns, and frequency
caps that guarantee no patient is contacted more than once a day
regardless of how many opportunity families matched them simultaneously.

\subsection{Expected-value ranking and causal uplift}

Rather than ranking by a single risk score, five trained models are
combined: \textbf{no-show risk} (will this patient miss the appointment
if booked), \textbf{booking propensity} (will they book at all),
\textbf{expected revenue} (what is this visit worth), \textbf{causal
uplift} (a T-learner estimate of the \emph{incremental} effect of
contacting this specific patient, as opposed to their baseline
probability of returning anyway), and \textbf{predicted time-to-book} (a
survival-analysis estimate of how long until they book). A patient
predicted to book anyway, with zero incremental uplift, should be ranked
low even if their raw booking-propensity score is high --- this is
precisely the distinction a naive single-model ranking cannot make.

\subsection{The causal holdout experiment}

A randomly-assigned slice of the population that is deliberately never
contacted, regardless of how highly they would otherwise rank. Comparing
the treated group's real booking rate against the holdout group's real
booking rate is the only way this system can claim its outreach
\emph{caused} incremental bookings rather than simply having reached
people who would have returned on their own.

\subsection{Clinical note intelligence}

Unstructured clinical notes pulled from Epic are processed by a medical
LLM (MedGemma) to extract structured diagnoses, medications, procedures,
and follow-up recommendations, gated by a content-safety model
(NemoGuard) and by non-LLM validation checks before anything is stored.
Only structured output is persisted; raw note text is never retained.

\subsection{MLOps and governance}

Every model is tracked, versioned, and promoted through MLflow. Rules,
guardrails, and the pipeline's daily orchestration are handled by
Temporal, giving every step of the daily run durable retries and a full
execution history --- the same infrastructure discipline the sister
no-show/slot-recovery track applies to Airflow-scheduled retraining,
applied here to a long-running, always-on workflow engine instead of a
batch scheduler.

\section{Requirement analysis}

\subsection{System actors}

\begin{table}[h]
\centering
\caption{System actors}
\begin{tabular}{@{}p{3.5cm}p{10cm}@{}}
\toprule
Actor & Role in the system \\
\midrule
Clinic staff & Views opportunities, campaigns, WhatsApp conversations, and
bookings; monitors day-to-day operation. \\
Clinic owner & Everything staff can see, plus revenue and the holdout
experiment's causal results. \\
Clinic administrator & Everything owner can see, plus model
quality/drift monitoring and pipeline operations health. \\
Patient & Receives outreach on WhatsApp; replies to book, confirm, or
opt out. \\
\bottomrule
\end{tabular}
\end{table}

\subsection{Functional requirements}

\begin{table}[h]
\centering
\caption{Functional requirements}
\begin{tabular}{@{}p{1.2cm}p{12cm}@{}}
\toprule
ID & Functional requirement \\
\midrule
FR1 & Detect every patient matching at least one of 16 opportunity-family
rules, once per clinic per day. \\
FR2 & Gate detected opportunities by per-class patient consent, contact
cooldown, and frequency cap. \\
FR3 & Collapse multiple simultaneous opportunity matches for one patient
into a single contact. \\
FR4 & Score every gated candidate with five trained ML models and
combine them into one ranking. \\
FR5 & Randomly assign a holdout slice that is never contacted,
independent of rank. \\
FR6 & Dispatch WhatsApp outreach to the highest-ranked, non-holdout
candidates, capped per day. \\
FR7 & Fall back to SMS automatically on WhatsApp delivery failure. \\
FR8 & Ingest real patient, condition, encounter, procedure, and coverage
data from Epic FHIR. \\
FR9 & Extract structured clinical information from unstructured notes,
gated by a safety check. \\
FR10 & Record delivery status, replies, bookings, attendance, and
revenue against each campaign. \\
FR11 & Allow manual outcome marking (booked/attended) when no automated
feed exists, writing to the same tables an automated event would. \\
FR12 & Compare the treated and holdout arms' real booking rates to
report a causal effect. \\
FR13 & Present all of the above through a role-gated web console
(staff/owner/admin tiers). \\
FR14 & Track every model's training runs, versions, and promotion
status. \\
\bottomrule
\end{tabular}
\end{table}

\subsection{Non-functional requirements}

\begin{table}[h]
\centering
\caption{Non-functional requirements}
\begin{tabular}{@{}p{3cm}p{10.5cm}@{}}
\toprule
Requirement & Description \\
\midrule
Safety & Holdout-arm patients can never be marked booked/attended
through the console --- doing so would corrupt the causal comparison the
holdout exists to protect. \\
Traceability & Every detected opportunity, gating decision, model score,
dispatch, and outcome is persisted and auditable. \\
Explainability & The no-show score is accompanied by a SHAP breakdown so
staff can see why a patient is considered high-risk. \\
Multi-tenancy & The platform serves multiple clinics from one deployment
with fully isolated per-clinic data. \\
Idempotency & Every write (opportunity, campaign, outcome) is
upsert-based; retrying an already-applied step never duplicates or
corrupts state. \\
Fault isolation & An Epic, WhatsApp, or clinical-note-extraction failure
degrades gracefully and never blocks the rest of the daily pipeline. \\
Consent-first & No outreach action is ever dispatched to a patient
without a matching, class-specific consent record. \\
\bottomrule
\end{tabular}
\end{table}

\subsection{Implemented scope}

\begin{table}[h]
\centering
\caption{Implemented and future scope}
\begin{tabular}{@{}p{2.6cm}p{11cm}@{}}
\toprule
Scope & Elements \\
\midrule
Implemented & Epic FHIR integration (OAuth interactive flow + Bulk Data
export), 16 YAML opportunity families, per-class consent gating,
cooldown/frequency caps, five trained ranking models plus a
lookalike-discovery model and a WhatsApp-reply intent classifier,
randomized holdout experiment, Temporal-orchestrated daily pipeline,
WhatsApp (Meta Cloud API) dispatch with Twilio SMS fallback,
clinical-note extraction (MedGemma + NemoGuard) on a GPU server,
MLflow-tracked model registry, a Jinja2 + React console with role-gated
pages, and full multi-clinic isolation. \\
Prototype / future work & A dental treatment-pathway state-machine model
(personalized next-clinical-step suggestion) --- deliberately scoped as
engineering prototype only, gated behind clinical-advisor review before
any real use (Section~\ref{sec:dental-prototype}). \\
Not implemented / external dependency & A real (non-sandbox) production
Epic connection for a live clinic, commercial-scale Meta WhatsApp
template approval, and cross-module integration with the company's
separate scheduling (Velo Desk) and billing/insurance (Velo Claim)
products. \\
\bottomrule
\end{tabular}
\end{table}

\section{Use case diagram}

\begin{figure}[h]
\centering
\fbox{\parbox{0.85\textwidth}{\centering\vspace{2cm}
TODO: insert the use case diagram here. \\
Two primary human actors --- \emph{Clinic Staff/Owner/Admin} and
\emph{Patient} --- mirror the sister report's diagram convention.
Clinic-side use cases: view opportunities, view campaigns and journey
status, review model scores and SHAP explanation, mark outcome
manually, view WhatsApp conversations, view revenue and holdout results
(owner), view model/pipeline health (admin). Patient-side use cases:
receive WhatsApp outreach, reply to book/confirm, opt out. Epic and
Meta/Twilio are external systems, not use-case actors, exactly as the
EHR is excluded from the sister report's diagram and explained instead
through the architecture chapter.
\vspace{2cm}}}
\caption{Use case diagram of the Velo Engage (VeloDoc) application}
\label{fig:usecase}
\end{figure}

\section{Project planning and development approach}

\subsection{Project team}

The project involved the same category of academic, industrial, and
technical guidance as the sister PFE track at Velodoc: an academic
supervisor for scientific orientation and report quality, industrial
supervisors (Mr.\ Wael Hilali, Mr.\ Bilel Said) for technical direction
and system validation, and the intern author responsible for design,
implementation, testing, evaluation, and documentation.

\subsection{Project development phases}

The work proceeded through recognizable phases even though it was run
iteratively rather than as a rigid waterfall: (1)~Epic integration and
data foundation; (2)~opportunity detection and guardrails; (3)~the five
ranking models and the holdout experiment; (4)~orchestration, dispatch,
and the console; (5)~clinical-note intelligence; (6)~UI modernization
and a full React migration; (7)~documentation and defense preparation.

\subsection{Iterative development approach}

Each subsystem was built, tested, and verified against the real running
stack before the next was layered on top --- Epic integration was
verified against a live sandbox patient before the ranking models were
trained against its output; the React console migration was verified
page-by-page against a real, signed session cookie hitting the live
container, confirming zero regression in the original server-rendered
pages before each new page was added.

\subsection{Planning timeline}

\begin{figure}[h]
\centering
\fbox{\parbox{0.85\textwidth}{\centering\vspace{1.5cm}
TODO: insert a Gantt-style planning figure, analogous to the sister
report's timeline figure, with phases: Epic integration $\rightarrow$
Opportunity detection \& guardrails $\rightarrow$ Ranking models \&
holdout $\rightarrow$ Orchestration, dispatch \& console $\rightarrow$
Clinical-note intelligence $\rightarrow$ UI modernization \& React
migration $\rightarrow$ Documentation \& defense prep.
\vspace{1.5cm}}}
\caption{Project planning timeline}
\label{fig:timeline}
\end{figure}

\section{Conclusion}

This chapter presented the internship context, the operational problem
of unmanaged patient re-engagement, the limits of existing recall
tooling, and the proposed Velo Engage solution: a rule-gated, ML-ranked,
causally-measured outreach pipeline. The next chapter presents the
technical architecture.

% =============================================================================
\chapter{System and AI Architecture}
\label{ch:architecture}

\section{Introduction}

This chapter presents the technical architecture: how the system is
organized around Epic as the clinical source of truth, the
Temporal-orchestrated daily pipeline, the five ranking models plus the
discovery and intent-classification models, the causal holdout
mechanism, and the MLOps governance design.

\section{Architectural objectives and constraints}
\label{sec:objectives}

\textbf{Epic remains the source of truth for clinical data.} Velo Engage
never invents or overrides clinical facts --- it reads patient,
condition, encounter, procedure, and coverage data from Epic's FHIR API,
and its own database stores only operational/predictive data
(opportunity records, model scores, campaigns, consent, outcomes).

\textbf{Rule-based safety, ML-based ranking --- never the reverse.}
Anything with a direct safety or compliance consequence (consent
gating, contact cooldowns, holdout-arm protection) is deterministic rule
logic. Machine learning is used exclusively to \emph{rank} or
\emph{score} within a space of options the rule layer has already deemed
valid --- never to decide validity itself. This single principle
governs every architectural decision described in this chapter and is
the same principle that, in this project's later prototype work
(Section~\ref{sec:dental-prototype}), explicitly excludes an LLM from
making any clinical-suggestion decision directly.

\textbf{Fault isolation per external dependency.} Epic, Meta/Twilio, and
the GPU-hosted clinical-note model are each isolated in their own
service, specifically so one external system's failure mode (an expired
Epic sandbox token, a WhatsApp delivery failure, a cold-starting GPU
model) cannot block the rest of the daily pipeline.

\textbf{Multi-clinic isolation.} Every table, every rule evaluation, and
every model score is scoped by \code{clinic_id}; the platform runs one
Temporal workflow per clinic per day, not one shared global run.

\textbf{Auditability of the causal claim.} Because the holdout
experiment is the system's central evidentiary mechanism, the
architecture treats ``never contact a holdout patient'' as an invariant
enforced at the console's write layer, not merely a convention observed
by the ranking logic.

\section{EHR-centered data architecture}

\code{ve_connect} is the only service that talks to Epic directly. It
implements two independent authentication flows with independent token
lifetimes: an \textbf{interactive OAuth (PKCE) flow} used for
individual-patient FHIR reads, and a separate \textbf{SMART Backend
Services (client-credentials/JWT) flow} used only for Bulk Data
\code{\$export} against a defined patient cohort (a Group). These two
flows can be in completely different states at once --- the interactive
token can be expired while Bulk Data exports keep succeeding --- because
they are unrelated credentials with unrelated lifetimes; this is a real,
recurring characteristic of testing against Epic's non-production
sandbox, not a defect.

\code{mapper.py} translates raw FHIR resources (\code{Patient},
\code{Condition}, \code{Encounter}, \code{Procedure}, \code{Coverage},
and --- for clinical-note intelligence --- \code{DocumentReference}/
\code{Binary}) into the fields the platform's own schema needs: ICD-10
codes, last visit date, last procedure, coverage end date, phone number.
Fetching is fault-tolerant per resource type --- one Epic resource type
being unavailable for this app's authorized scope degrades one derived
feature to a safe default rather than blocking every other resource type
Epic does authorize.

A diagnostic-only endpoint exposes the \emph{exact} raw FHIR resources
the real pipeline sees for a given patient, before any mapping --- used
to confirm data-quality issues (e.g., sandbox test/filler clinical
notes) are a property of the sandbox data, not a pipeline defect,
without ever persisting the raw payload to disk.

\begin{figure}[h]
\centering
\fbox{\parbox{0.85\textwidth}{\centering\vspace{2cm}
TODO: reproduce the system architecture diagram from
\texttt{ARCHITECTURE.md} (the repository's own Mermaid diagram) here as
a proper figure: External (Epic, Meta, Twilio) $\rightarrow$ Core
Pipeline (\texttt{ve_connect}, \texttt{ve_orchestrator}, \texttt{ve_reach},
\texttt{ve_measure}) $\rightarrow$ Data (Postgres, Neo4j, Redis, MLflow)
$\rightarrow$ Intelligence (5 ML models, \texttt{ve_clinical_intel}) and
\texttt{ve_console} on top.
\vspace{2cm}}}
\caption{Global architecture of the Velo Engage (VeloDoc) application}
\label{fig:architecture}
\end{figure}

\section{Global system architecture}

The platform is organized as a set of small, single-purpose services
coordinated by Temporal and sharing a common Postgres database as the
operational source of truth:

\begin{itemize}
  \item \textbf{\code{ve_connect}} --- Epic FHIR integration (patient
        data and clinical notes).
  \item \textbf{\code{ve_orchestrator}} --- the Temporal-hosted daily
        workflow: detect, gate, rank, dispatch.
  \item \textbf{\code{ve_reach}} --- WhatsApp/SMS sending and inbound
        webhook handling (delivery status, replies, opt-outs).
  \item \textbf{\code{ve_measure}} --- booking and revenue attribution
        from real-world events.
  \item \textbf{\code{ve_console}} --- the web application (Jinja2
        server-rendered pages and a parallel React single-page
        application, both serving identical data through a shared JSON
        API layer).
  \item \textbf{\code{ve_clinical_intel}} --- clinical-note structured
        extraction, running on a GPU server (DGX Spark) separate from
        the rest of the stack.
\end{itemize}

Supporting infrastructure: \textbf{PostgreSQL} (patients, campaigns,
outcomes, consent, revenue --- the shared operational store nearly
every service reads/writes), \textbf{Neo4j} (the clinical-recall
knowledge graph queried by one opportunity family), \textbf{Redis} (a
short-lived feature cache in front of Postgres during ranking),
\textbf{MLflow} (every model's training runs, versions, and the
production pointer the daily pipeline loads at ranking time),
\textbf{Keycloak} (OIDC login for the console), \textbf{HashiCorp Vault}
(every service's credentials, read once at startup with an
environment-variable fallback so a transient Vault outage never blocks
the pipeline), and \textbf{NATS JetStream} (the one genuinely
asynchronous handoff in the system: \code{ve_orchestrator} publishes a
dispatch-ready campaign, \code{ve_reach} consumes and sends it,
decoupled so a slow or unavailable send never blocks the ranking
pipeline itself).

Communication patterns follow a consistent rule: Postgres for shared
state that any service can read directly; NATS JetStream for the one
asynchronous cross-service handoff; plain HTTP for direct
request/response between services with no shared database access
(\code{ve_clinical_intel}, on the GPU server, calling \code{ve_connect}'s
REST API); and inbound webhooks for external events (Meta delivering
delivery-status and reply events to \code{ve_reach}). Epic never pushes
anything --- \code{ve_connect} pulls from Epic on the orchestrator's own
schedule.

\section{Main pipeline workflows}

\subsection{The daily detect--gate--rank--dispatch--measure workflow}
\label{sec:dailyworkflow}

Every service above is orchestrated by a single Temporal workflow
(\code{DailyEngagementWorkflow}), running once per clinic per day,
composed of retryable, independently-executed activities:

\begin{enumerate}
  \item \textbf{Refresh data} --- a best-effort pull of the latest Epic
        data and a recomputation of behavioral features (visit history,
        no-show history, message engagement). This activity
        deliberately swallows its own errors: an Epic hiccup should
        degrade gracefully, not halt the pipeline.
  \item \textbf{Detect} --- the 16 opportunity-family rules are
        evaluated against every patient's current feature row.
  \item \textbf{Gate} --- candidates are filtered to those with matching
        per-class consent, respecting cooldown and frequency-cap
        constraints, and collapsed to one contact per patient per day.
  \item \textbf{Rank} --- the five ranking models score every survivor
        and combine into one expected-value ranking; a randomly-selected
        slice is independently assigned to the holdout arm regardless of
        rank.
  \item \textbf{Dispatch} --- the highest-ranked, non-holdout candidates
        (a capped daily number) are sent to \code{ve_reach} for WhatsApp
        delivery, with automatic SMS fallback.
  \item \textbf{Measure} --- delivery status, replies, bookings,
        attendance, and revenue flow back automatically via webhooks (or
        are recorded manually where no automated feed exists), closing
        the loop for the next day's ranking and eventual model
        retraining, and producing the ROI figure shown on the console's
        dashboard.
\end{enumerate}

\begin{figure}[h]
\centering
\fbox{\parbox{0.85\textwidth}{\centering\vspace{1.5cm}
TODO: render this six-step pipeline as a flowchart figure, mirroring
Figure~6 in the sister report's style.
\vspace{1.5cm}}}
\caption{The daily detect--gate--rank--dispatch--measure workflow}
\label{fig:pipeline}
\end{figure}

\subsection{Opportunity detection}

Sixteen independent, YAML-defined ``opportunity families'' --- most
deterministic threshold rules (dormancy days, benefit-expiry windows,
open treatment-plan flags), one backed by a Neo4j clinical-guideline
graph, and one backed by a lookalike/collaborative-filtering model that
recommends a family a patient has never used based on similar patients
who have.

\subsection{Guardrails and consent}

Per-class consent checking (a patient consenting to a routine recall
does not imply consent to a promotional contact), contact cooldowns,
frequency caps, and de-duplication so a patient matching multiple
opportunity families in the same run is still contacted at most once.

\subsection{The causal holdout experiment}

A randomly-assigned slice of the ranked population that is never
contacted, entirely independent of how highly it would otherwise rank.
Comparing this arm's real booking rate against the treated arm's is the
platform's mechanism for proving incremental causal impact rather than
correlation --- and this arm's protection is enforced structurally: the
console's outcome-marking endpoint rejects, with an explicit
\code{400} response and explanation shown in the UI, any attempt to mark
a holdout patient as booked or attended, since doing so would corrupt the
exact comparison the holdout exists to protect.

Formally, let $T_i \in \{0,1\}$ indicate whether patient $i$ was
assigned to the treated (contacted) arm, and let $Y_i \in \{0,1\}$
indicate whether patient $i$ booked. Randomization of $T_i$ (independent
of the ranking score) allows the platform to estimate the
\textbf{average treatment effect} of outreach as a simple difference in
means between arms:

\begin{equation}
\widehat{\mathrm{ATE}} = \frac{1}{n_1}\sum_{i:\,T_i=1} Y_i \;-\;
\frac{1}{n_0}\sum_{i:\,T_i=0} Y_i
\label{eq:ate}
\end{equation}

where $n_1$ and $n_0$ are the number of patients in the treated and
holdout arms respectively. This is the same quantity the causal uplift
model (Section~\ref{sec:uplift}) is trained to \emph{predict}
per-patient, before any outcome is observed.

\subsection{WhatsApp dispatch and reply handling}

WhatsApp (Meta Cloud API) is the primary channel; Twilio SMS is the
automatic fallback on delivery failure. Inbound replies are parsed for
booking intent and opt-out requests by a trained intent classifier, and
every message (outbound and inbound) is logged for the console's
conversation view.

\subsection{Clinical note intelligence}

Unstructured clinical notes pulled from Epic (\code{DocumentReference}/
\code{Binary}) are passed through a content-safety check (NemoGuard)
before ever reaching the extraction model; a medical LLM (MedGemma) then
extracts structured diagnoses, medications, procedures, and follow-up
recommendations; non-LLM validation checks (well-formed code shapes,
non-empty fields) flag anything that doesn't parse cleanly rather than
silently discarding it. Only the structured output is persisted --- raw
note text is discarded after extraction, matching this project's
no-raw-retention discipline for the most sensitive category of data it
ever touches.

\section{Model architecture and technical choices}

\subsection{No-show risk model}

A CatBoost/LightGBM gradient-boosted classifier trained on a
research-guided synthetic dataset of 110{,}000 UAE/KSA-style appointment
records, using booking-time and historical behavioral features (prior
no-show ratio, consecutive no-show streak, lead time, same-day booking
flag, age, deposit status, travel time, cancellation history, VIP
status, new-patient flag, urgency, neighborhood no-show rate, and
provider effect, among others). As in every gradient-boosted-ensemble
model in this platform, the model score is an additive ensemble of
regression trees mapped through a logistic link:

\begin{equation}
s(x) = b + \sum_{k=1}^{K} \eta\, T_k(x), \qquad
r_0 = P(y=1\mid x) = \sigma(s(x)) = \frac{1}{1+\exp(-s(x))}
\label{eq:catboost}
\end{equation}

where $x$ is the feature vector, $b$ is the learned bias, $T_k(x)$ is
the output of tree $k$, $K$ is the number of trees, $\eta$ is the
learning rate, and $\sigma$ is the sigmoid function. Gradient-boosted
trees were chosen over deep-learning alternatives for the same reason
this choice is made throughout the platform's other tabular scoring
tasks: strong performance on mixed categorical/numerical tabular data
without heavy manual encoding, native SHAP explainability, and
millisecond-scale inference cost appropriate for scoring an entire
clinic's patient population daily.

\subsection{Booking propensity model}

Estimates the probability a candidate books at all if contacted,
independent of no-show risk --- a patient can be very likely to book
yet, once booked, still carry meaningful no-show risk; conflating the
two into one score would lose exactly the distinction the ranking
formula needs.

\subsection{Expected-value (revenue) model}

Estimates the expected monetary value of a visit for a given
patient/opportunity combination, so the ranking can prefer a
lower-propensity but much higher-value opportunity over a
high-propensity, low-value one when that trade-off is worth it.

\subsection{Causal uplift model (T-learner)}
\label{sec:uplift}

A two-model (treatment/control) T-learner estimating the
\emph{incremental} effect of contact for a specific patient. Two models
are fit separately --- $\mu_1(x)$ on the historically-treated population
and $\mu_0(x)$ on the historically-untreated (or holdout) population ---
and the uplift estimate is their difference:

\begin{equation}
\hat\tau(x) = \hat\mu_1(x) - \hat\mu_0(x)
\label{eq:uplift}
\end{equation}

This is why the uplift score can legitimately be negative: for a patient
predicted to book slightly \emph{less} often when contacted than when
left alone (a plausible real phenomenon --- e.g., a contact that reads
as impersonal spam to a patient who would have called in on their own),
the honest estimate from Equation~\ref{eq:uplift} is negative, and the
ranking formula applies a steep discount to negative-uplift candidates
rather than clipping the score to zero and hiding the signal.

\subsection{Survival model (predicted time-to-book)}

A time-to-event model estimating how long until a given patient books if
contacted, letting the platform distinguish ``will book soon'' from
``will book eventually'' opportunities when the daily dispatch cap
forces a choice between them.

\subsection{Lookalike/discovery model}

Used only for one opportunity family: a collaborative-filtering-style
model that recommends a family a patient has never used, based on the
behavior of similar patients --- the one place in opportunity
\emph{detection} (as opposed to ranking) where a model, rather than a
deterministic rule, decides whether a candidate reason exists at all,
since ``similarity to other patients'' is inherently a learned
relationship rather than a fixed threshold.

\subsection{WhatsApp reply intent classifier}

A lightweight classifier trained to distinguish booking intent,
confirmation, cancellation/opt-out, and ambiguous replies from inbound
WhatsApp messages, feeding the reply-handling logic in \code{ve_reach}.

\subsection{The combined ranking score}

The five models' outputs are combined into a single expected-value
score used to order the day's dispatch candidates:

\begin{equation}
\mathrm{EV}(x) = \hat p_{\mathrm{book}}(x)\cdot\big(1-\hat r_{\mathrm{noshow}}(x)\big)\cdot \hat v(x) \cdot g\big(\hat\tau(x)\big)
\label{eq:ranking}
\end{equation}

where $\hat p_{\mathrm{book}}(x)$ is the booking-propensity score,
$\hat r_{\mathrm{noshow}}(x)$ is the no-show risk score, $\hat v(x)$ is
the expected-revenue score, and $g(\cdot)$ is a monotone discount
function of the uplift estimate $\hat\tau(x)$ that penalizes
negative-uplift candidates steeply rather than merely zeroing them out
--- consistent with the discussion in Section~\ref{sec:uplift}. The
predicted time-to-book score is used as a secondary tie-breaker within
the daily dispatch cap rather than as a multiplicative factor, since it
answers a different operational question (\emph{when}, not
\emph{whether}) from the other four models.

\subsection{Why not a single end-to-end model, and why not an LLM for ranking}

A single model predicting ``will this patient book'' cannot, by
construction, separate \emph{causation} from \emph{correlation} --- this
is precisely the gap the dedicated uplift model and the randomized
holdout close together. Equally, an LLM was deliberately not used
anywhere in the ranking or detection decision itself: ranking is a
structured, tabular-feature scoring problem gradient-boosted trees solve
well, cheaply, and with a clean SHAP-based explanation; an LLM applied
to that problem would add hallucination risk, latency, and cost while
solving nothing an LLM is actually good at. The one place a language
model is used at all --- clinical-note extraction --- is precisely the
one place the problem genuinely requires understanding unstructured
text, and even there it is bounded by a safety-content gate and non-LLM
validation before its output is trusted.

\section{MLOps and governance architecture}

Every model is trained via a dedicated script (\code{ml/training/train_*.py}),
tracked in \textbf{MLflow} with its metrics, and served in production via
a paired registry/scorer module (\code{ml/registry/*_scorer.py}) that the
daily Temporal workflow loads at ranking time. Promotion from a
newly-trained candidate to the production pointer
(\code{ml/registry/promote.py}) is a deliberate, auditable step rather
than automatic --- mirroring, in spirit, the same ``candidate must clear
evaluation and audit gates before the runtime system uses it''
discipline used elsewhere in the Velodoc team's MLOps designs, adapted
here to a registry-pointer promotion rather than a scheduled retraining
DAG, since the ranking models are retrained on a slower, less
time-critical cadence than a per-message online scorer would need.

\begin{figure}[h]
\centering
\fbox{\parbox{0.85\textwidth}{\centering\vspace{1.5cm}
TODO: render the MLOps loop as a figure: training script $\rightarrow$
MLflow tracking $\rightarrow$ evaluation/audit $\rightarrow$ registry
promotion $\rightarrow$ daily workflow load, mirroring the sister
report's MLOps-loop figure style.
\vspace{1.5cm}}}
\caption{MLOps and model-registry loop}
\label{fig:mlops}
\end{figure}

\section{Conclusion}

This chapter presented the EHR-centered data architecture, the global
service architecture, the daily detect--gate--rank--dispatch--measure
pipeline, the five ranking models plus the discovery and intent models,
and the MLOps governance design. The next chapter describes how this
architecture was implemented.

% =============================================================================
\chapter{Implementation of the Velo Engage (VeloDoc) Application}
\label{ch:implementation}

\section{Introduction}

This chapter describes how the architecture in Chapter~\ref{ch:architecture}
was translated into a working system: the technology stack, and the
implementation of each service in the pipeline, closing with an honest
account of a scoped, deliberately-limited prototype extension explored
beyond the core mission.

\section{Development environment and technology stack}

\begin{table}[h]
\centering
\caption{Main development and execution tools}
\begin{tabular}{@{}p{4cm}p{9.5cm}@{}}
\toprule
Tool / component & Role \\
\midrule
Python & Primary language for every backend service, training script, and
workflow activity. \\
FastAPI & API framework for every service (\code{ve_connect},
\code{ve_orchestrator}'s HTTP surfaces, \code{ve_reach}, \code{ve_measure},
\code{ve_console}, \code{ve_clinical_intel}). \\
Temporal & Workflow engine orchestrating the daily
detect--gate--rank--dispatch pipeline, one workflow per clinic per day,
with durable retries and full execution history. \\
PostgreSQL & Shared operational store: patients, campaigns, outcomes,
consent, revenue. \\
Neo4j & The clinical-recall knowledge graph queried by one opportunity
family. \\
Redis & Short-lived feature cache in front of Postgres during ranking. \\
MLflow & Experiment tracking and model registry for all trained models. \\
CatBoost / LightGBM / scikit-learn & The ranking, discovery, and
intent-classification models. \\
MedGemma + NVIDIA NemoGuard & Clinical-note structured extraction and
its content-safety gate, hosted on a DGX Spark GPU server. \\
Meta WhatsApp Business Cloud API, Twilio & Primary outreach channel and
its automatic SMS fallback. \\
Epic FHIR (SMART on FHIR, Bulk Data Export) & Real patient data and
clinical-note ingestion. \\
Jinja2 + React/TypeScript (Vite) & The console's two frontends, serving
identical data through a shared JSON API. \\
Keycloak (OIDC) & Console authentication. \\
HashiCorp Vault & Centralized secrets, with an environment-variable
fallback if unreachable. \\
NATS JetStream & The asynchronous orchestrator$\rightarrow$reach
dispatch handoff. \\
Docker / Docker Compose & Reproducible multi-service local and DGX
Spark deployment. \\
\bottomrule
\end{tabular}
\end{table}

\section{Epic integration implementation}

\code{ve_connect} implements \code{auth.py} (interactive PKCE OAuth),
\code{bulk_auth.py} (SMART Backend Services JWT client-assertion for Bulk
Data), \code{fhir_client.py} (resource fetch/search with automatic
Bundle pagination), \code{epic_bulk.py} (the \code{\$export} flow
against a Group, polling its status URL, downloading the resulting
NDJSON files, and deleting the export afterward), \code{mapper.py}
(FHIR-to-internal-schema translation), and \code{adapter.py} (the write
path into \code{staging_patients}/\code{patient_features}/\code{consent},
upserted in one transaction per patient). Two real, silent bugs were
found and fixed during Epic sandbox testing: \code{Encounter.class}
filtering initially assumed standard FHIR coding and matched zero real
Epic encounters, since Epic populates it with a proprietary internal
code system; and \code{CarePlan} fetching, rejected by the app's
authorized scope in two different ways, was made fault-tolerant so one
unavailable resource type degrades a single derived flag rather than
blocking every other resource type the app \emph{is} authorized for.

\section{Orchestration implementation}

\code{ve_orchestrator} implements \code{workflows.py} (the
\code{DailyEngagementWorkflow} Temporal definition), \code{activities.py}
(each pipeline stage as an independently retryable activity ---
activities pass only IDs between each other and re-fetch their own data
from Postgres, keeping every activity result well under Temporal's
payload limits), \code{policy_engine.py} (loading and safely evaluating
the YAML opportunity-family rules, with a validator that rejects a
malformed policy file at load time rather than at runtime),
\code{graph_client.py} (the Neo4j clinical-recall query),
\code{feature_store.py} (the Redis-backed feature cache), and
\code{historical_features.py} (a dedicated ETL recomputing the no-show
model's required historical features, since patient behavioral history
is not something Epic itself exposes as a ready-made feature set).

\section{Database design}

A single PostgreSQL instance holds every operational table, with
\code{clinic_id} on every row as the multi-tenant isolation key.
Migrations are applied in strict numerical order
(\code{infra/migrations/001_*.sql} through the current head), each
migration documenting, in its own header comment, exactly what gap it
closes and why --- a discipline carried through every schema change made
across this project, including the dental-pathway prototype's own
migration (Section~\ref{sec:dental-prototype}).

\section{Opportunity detection and guardrails implementation}

Opportunity families are plain YAML files
(\code{policies/opportunities/*.yml}), each declaring a \code{family} id,
a \code{consent_class}, a \code{priority}, and a set of conditions
evaluated against a patient's feature row by pure, independently
unit-tested functions (\code{conditions_match}, \code{evaluate_policy},
\code{build_evidence}) --- deliberately kept free of any database or
Temporal dependency so the rule-matching logic itself can be tested
exhaustively without standing up any infrastructure. Guardrails
(consent-class checking, cooldown, frequency cap, de-duplication) are
implemented as a separate gating activity applied after detection and
before ranking.

\section{Ranking model implementation}

Each of the five models follows the same paired convention: a training
script (\code{ml/training/train_*.py}) that produces an MLflow-tracked
artifact, and a registry/scorer module (\code{ml/registry/*_scorer.py})
with a documented \code{DEFAULTS} fallback used whenever a patient's
feature row is incomplete --- so a missing historical feature degrades
to a known, safe default score rather than raising an exception or
silently returning a wrong ranking. The no-show model's 32-feature
contract is enforced end to end: a dedicated migration and ETL
(\code{historical_features.py}) exist specifically to keep
\code{patient_features} populated with every column the scorer's
feature contract requires, closing a gap where several columns
previously fell back to their \code{DEFAULTS} value for every patient
regardless of that patient's real history.

\section{Causal holdout implementation}

Holdout assignment happens inside the same ranking activity that
computes the expected-value score, using a deterministic, seeded random
assignment so a given patient's holdout status is reproducible for a
given run rather than re-rolled arbitrarily. The console enforces the
holdout arm's protection at the write layer: the outcome-marking
endpoint checks a campaign's arm before accepting a ``booked'' or
``attended'' mark, returning an explicit \code{400} with an explanation
surfaced in the UI (rather than the buttons simply not appearing) if a
holdout campaign is targeted, so a coordinator understands \emph{why}
the action is refused rather than assuming a bug.

\section{Communication channel implementation}

\code{ve_reach} sends WhatsApp messages via the Meta Cloud API and
automatically falls back to Twilio SMS on delivery failure; inbound
webhooks from Meta (delivery status, replies) are received and routed to
the intent classifier for reply parsing. A real deployment constraint
surfaced during this work: ngrok's free tier grants exactly one reserved
public domain, which both WhatsApp's inbound webhook and Epic's OAuth
redirect needed simultaneously. This was solved permanently, not with a
manual workaround, by adding three proxy routes to \code{ve_reach}
(\code{/auth/login}, \code{/auth/callback}, \code{/}) that forward
Epic's OAuth traffic over the internal Docker network to
\code{ve_connect}'s real handlers and relay the response verbatim --- one
public tunnel now serves both purposes permanently, verified live
through the actual public tunnel rather than only against localhost.

\section{Clinical note intelligence implementation}

\code{ve_clinical_intel} runs on a separate DGX Spark GPU server,
calling \code{ve_connect}'s REST API to fetch notes (never given direct
database access from the GPU host) and to store extractions. The
extraction pipeline is: content-safety check (NemoGuard) first --- a
note that fails this check is never sent to the extraction model at all
--- then structured extraction (MedGemma, a 27-billion-parameter
text-focused medical model, accessed via a streaming server-sent-events
API accumulated into a full response before parsing), then non-LLM
validation of the parsed structure, then storage. A real, non-obvious
finding from this work: some real Epic sandbox patients' clinical notes
are themselves sandbox connectivity-test filler text or
repeated-character load-test artifacts, not genuine clinical
documentation --- explaining, correctly, why extraction sometimes
returns an empty result for a given note. This was confirmed, not
assumed, by inspecting the actual raw note text through a diagnostic
endpoint that returns exactly what the extraction pipeline itself
receives, before concluding anything was broken in the extraction logic.

\section{Console implementation}

\code{ve_console} is implemented twice over the same data and the same
authentication: a server-rendered Jinja2 application (the original
implementation) and a React + TypeScript single-page application mounted
at \code{/app}, both calling the \emph{same} underlying data-fetch
functions (\code{_overview_data()}, \code{_campaigns_data()}, and so on
for every page) so the two frontends can never drift apart in what they
display. The React migration was executed incrementally, one page
verified before the next, using a genuinely-signed session cookie to
test the real running container end to end --- confirming every original
page's behavior was unchanged \emph{after} the migration, not merely
assumed unchanged before it. Role gating (staff/owner/admin) is
additive, checked by testing for either role granting a given tier
rather than a separate expansion step.

\begin{figure}[h]
\centering
\fbox{\parbox{0.85\textwidth}{\centering\vspace{2cm}
TODO: insert a screenshot of the console overview/campaign-detail page,
including the model-scores panel (expected value, no-show risk, booking
propensity, expected revenue, uplift, predicted time-to-book) with its
SHAP explanation panel.
\vspace{2cm}}}
\caption{Console campaign-detail view showing the model scores and SHAP explanation}
\label{fig:console}
\end{figure}

\section{MLOps implementation}

Model training scripts, MLflow experiment tracking, and the
registry-promotion step (\code{ml/registry/promote.py}) together form the
platform's MLOps surface. Every registered model version is queryable
from the console's model-health page (role-gated to administrators),
alongside drift and quality metrics for the no-show model specifically,
since it is the model with the most mature retraining pipeline in this
implementation.

\section{Dockerization and deployment}

The full stack (Postgres, Neo4j, Redis, NATS, Temporal, Keycloak, Vault,
MLflow, and every service above) is defined in a single
\code{docker-compose.dev.yml}, with \code{ve_console}'s image built as a
two-stage Docker build (a Node.js stage compiling the React frontend,
copied into the final Python runtime image so the running container
never needs Node at all). \code{ve_clinical_intel} and the underlying
model server run on the DGX Spark GPU host, reachable from the rest of
the stack over the team's shared network.

\section{Prototype extension: personalized dental treatment-pathway suggestion}
\label{sec:dental-prototype}

Beyond the core reactivation-and-measurement mission, an exploratory
extension was designed and partially scaffolded: a state-machine model
that suggests the clinically appropriate \emph{next step} in a dental
treatment sequence (e.g., a confirmed deep-caries diagnosis
$\rightarrow$ root canal $\rightarrow$ crown), personalized per patient
rather than looked up from a generic protocol.

This is deliberately treated as a \textbf{prototype, not a production
feature}, for reasons consistent with this whole project's governing
principle (Section~\ref{sec:objectives}): a wrong clinical suggestion
carries real consequence, so the design places a clinician-authored,
YAML-defined rule graph as the sole authority over which transitions are
\emph{valid}, with machine learning used only to \emph{rank} among the
options the graph already permits --- never to invent a transition. Each
patient-case's position in the graph is modeled as a Temporal workflow
instance (Temporal workflows already are durable state machines, so this
reuses the same orchestration engine the core pipeline runs on rather
than building a second one), advanced by a signal when new procedure
data arrives and resolved by a clinician's explicit accept/override
decision whenever more than one transition is simultaneously valid. A
sibling service (\code{ve_connect_dental}) was scaffolded to pull real
dental procedure history from OpenDental (a practice-management system
distinct from Epic, since standalone dental clinics typically do not run
Epic), behind a transport-agnostic client interface so a
direct-database Phase~A (self-hosted, requiring no external
registration) can later be swapped for OpenDental's official REST API
without touching the mapping or state-machine logic above it.

Nothing produced by this prototype has been reviewed by a dental
clinical advisor, and the codebase enforces this honestly rather than
silently: every new file in this extension states its \code{prototype}
status explicitly, and the design's own documentation records, as an
open item, that production use requires exactly two things this
internship's scope did not include --- a named clinical reviewer signing
off on the pathway graph, and a confirmed real dental-PMS data source
for a specific pilot clinic. It is presented here only as evidence of
the architecture's extensibility, not as delivered scope.

\section{Conclusion}

This chapter described the implementation of every service in the
pipeline: Epic integration, orchestration, opportunity detection and
guardrails, the five ranking models, the causal holdout, WhatsApp/SMS
dispatch, clinical-note intelligence, the console, MLOps, deployment,
and an honestly-scoped prototype extension. The next chapter presents
the evaluation and validation of the implemented system.

% =============================================================================
\chapter{Evaluation and Validation}
\label{ch:evaluation}

\section{Introduction}

This chapter presents the evaluation methodology, the results actually
obtained, and this project's honest limitations. Consistent with the
rest of this report, every number cited below is drawn directly from
this project's own repository artifacts (\code{ml/data/},
\code{ml/data/reports/}) rather than restated from memory, and every
cell that depends on a metric not yet re-extracted at the time of
writing is marked for the author to fill from the current MLflow
registry before submission, rather than filled with an invented figure.

\section{Evaluation methodology}

The no-show model was trained and evaluated on a research-guided
synthetic dataset of \textbf{110{,}000 UAE/KSA-style appointment
records} (file: \code{synthetic_no_show_uae_ksa_110k}), constructed the
same way the sister PFE track's own no-show simulator was built ---
from probability effects extracted from outpatient no-show literature
rather than arbitrary random labels --- so the simulated label follows a
realistic risk structure. The dataset's observed base no-show rate is
\textbf{22.0\%} (mean of the \code{no_show} column across all rows),
with feature distributions summarized directly from the dataset
(Table~\ref{tab:features}). Evaluation used held-out test splits and
standard probabilistic-classification metrics (ROC-AUC, PR-AUC, Brier
score, expected calibration error, precision/recall/F1), consistent
with the rest of this platform's scoring tasks being trained on tabular
gradient-boosted models.

\begin{table}[h]
\centering
\caption{Selected feature distributions, 110k synthetic no-show dataset}
\label{tab:features}
\begin{tabular}{@{}lrrrr@{}}
\toprule
Feature & Mean & Std & Min & Max \\
\midrule
prior_no_show_ratio & 0.173 & 0.143 & 0.0 & 0.903 \\
consecutive_no_show_streak & 0.189 & 0.487 & 0.0 & 6.0 \\
lead_time_days & 11.54 & 9.23 & 0.0 & 87.0 \\
same_day_flag & 0.102 & 0.302 & 0.0 & 1.0 \\
age & 36.49 & 12.09 & 18.0 & 80.0 \\
deposit_paid & 0.383 & 0.486 & 0.0 & 1.0 \\
travel_time_minutes & 22.69 & 9.40 & 3.0 & 63.4 \\
cancellation_history_ratio & 0.143 & 0.088 & 0.0003 & 0.589 \\
vip_status & 0.090 & 0.286 & 0.0 & 1.0 \\
new_patient & 0.280 & 0.449 & 0.0 & 1.0 \\
urgency_score & 5.68 & 2.31 & 1.0 & 10.0 \\
neighbourhood_ns_rate & 0.092 & 0.038 & 0.050 & 0.226 \\
no_show (target) & 0.220 & 0.414 & 0.0 & 1.0 \\
\bottomrule
\end{tabular}
\end{table}

\begin{table}[h]
\centering
\caption{Top feature-importance contributors (no-show model)}
\label{tab:importance}
\begin{tabular}{@{}lr@{}}
\toprule
Feature & Relative importance (\%) \\
\midrule
same_day_flag & 19.62 \\
lead_time_days & 11.77 \\
specialty & 10.62 \\
prior_no_show_ratio & 9.57 \\
confirmation_status & 6.50 \\
cancellation_history_ratio & 5.40 \\
treatment_stage & 5.39 \\
deposit_paid & 4.33 \\
travel_time_minutes & 4.12 \\
whatsapp_engagement & 3.15 \\
\bottomrule
\end{tabular}
\end{table}

\section{No-show risk model evaluation}

The no-show model was evaluated across two staged configurations
reflecting how much lead time is available before an appointment: a
\textbf{booking-time} stage (scored the moment an appointment is
created) and a \textbf{pre-appointment} stage (scored closer to the
visit, with more accumulated behavioral signal available).
Table~\ref{tab:twostage} reports the two-stage evaluation actually
recorded for this model.

\begin{table}[h]
\centering
\caption{No-show model, two-stage evaluation}
\label{tab:twostage}
\begin{tabular}{@{}lrrrr@{}}
\toprule
Stage & Threshold & ROC-AUC & PR-AUC & Brier score \\
\midrule
Booking-time & 0.250 & 0.7353 & 0.4900 & 0.1471 \\
Pre-appointment & 0.270 & 0.7514 & 0.5336 & 0.1417 \\
\bottomrule
\end{tabular}
\end{table}

As expected, the pre-appointment stage --- scored with more accumulated
signal --- improves both discrimination (ROC-AUC) and probability
quality (Brier score) over the booking-time stage, confirming the value
of the online/behavioral feature refresh the pipeline performs before
each day's ranking run (Section~\ref{sec:dailyworkflow}).

A second evaluation compared two candidate feature-encoding strategies
for handling missing historical data --- a ``two-stage'' approach and a
``single-missingness-indicator'' approach --- across both the booking
and pre-appointment stages, on both train and held-out test splits
(Table~\ref{tab:encoding}), to decide which encoding strategy to
standardize on.

\begin{table}[h]
\centering
\caption{Feature-encoding strategy comparison (test split)}
\label{tab:encoding}
\begin{tabular}{@{}llrrr@{}}
\toprule
Approach & Stage & ROC-AUC & PR-AUC & Brier \\
\midrule
Two-stage encoding & Booking & 0.6695 & 0.4552 & 0.1838 \\
Single-missingness indicator & Booking & 0.6702 & 0.4565 & 0.1836 \\
Two-stage encoding & Pre-appointment & 0.7333 & 0.5420 & 0.1698 \\
Single-missingness indicator & Pre-appointment & 0.7327 & 0.5419 & 0.1699 \\
\bottomrule
\end{tabular}
\end{table}

The two strategies perform nearly identically, which is itself a useful
negative result: it means the simpler, more maintainable encoding
strategy can be adopted without sacrificing measurable model quality ---
consistent with this project's general preference for the simplest
design that meets the requirement, rather than the most sophisticated
one.

\section{Other ranking models (booking propensity, expected value, uplift, survival) and the discovery/intent models}

TODO: fill this section from your current MLflow registry before
submission --- do not restate a number from an earlier draft or from
memory. For each model, report: the metric(s) appropriate to its task
(ROC-AUC/PR-AUC for propensity; MAE/RMSE for expected-value regression;
the Qini coefficient or a similar uplift-specific metric for the
T-learner uplift model, since ROC-AUC is not a meaningful metric for
uplift; concordance index for the survival model; top-$k$ precision or
silhouette-style similarity metrics for the lookalike model;
accuracy/F1 for the WhatsApp intent classifier), the dataset it was
evaluated on, and the currently-registered MLflow version. This section
is intentionally left as a template rather than populated with
placeholder numbers, because --- unlike the no-show model above --- a
currently-recorded evaluation artifact for these five models could not
be verified at the time of writing; stating a number here without that
verification would misrepresent the project's actual results.

\section{Guardrail and consent verification}

The guardrail and consent-gating logic (per-class consent matching,
cooldown, frequency cap, de-duplication) is covered by dedicated unit
tests exercising the pure rule-matching functions directly --- asserting,
among other cases, that a missing feature causes a rule to correctly
evaluate as non-matching rather than raising an exception, that a
\code{contains_prefix} condition correctly matches both bare and
full-precision ICD-10 codes (since synthetic seed data and real Epic
data represent the same clinical category at different precision), and
that nested \code{all}/\code{any} condition trees evaluate correctly.
The holdout-arm protection was verified functionally against the live
console: an attempt to mark a holdout campaign's outcome as booked or
attended is rejected with an explicit \code{400} response, and the
corresponding UI buttons are replaced with an explanation rather than
silently hidden.

\section{Console and React-migration verification}

The console's dual-frontend architecture (Jinja2 and React, sharing one
data-fetch layer and one JSON API) was verified end to end using a
genuinely-signed session cookie --- built to match the authentication
middleware's exact cookie format --- to exercise the real running
container rather than a bypassed test harness. This confirmed: every one
of the eleven migrated pages' new JSON API endpoints return real, live
data consistent with the original server-rendered page (including
nested cases such as the models page's full set of registered model
cards and versions, and the holdout page's real per-arm booking rates);
every React route loads correctly on a direct page load as well as
through client-side navigation; the outcome-marking mutation writes
correctly to both the bookings and outcomes tables and correctly rejects
a holdout-campaign attempt through both the HTML form and the JSON API
path; and --- the specific regression risk a UI rewrite of this kind
carries --- every original Jinja route continues to return identical
output after the migration, not merely before it.

\section{End-to-end system testing}

End-to-end verification exercised the full pipeline against a real Epic
sandbox connection rather than only synthetic data: a live patient
cohort was pulled through the interactive OAuth flow, opportunity
families were evaluated against the resulting feature rows, guardrails
were applied, the five models scored the surviving candidates, holdout
assignment was applied, and a real WhatsApp message was dispatched and
its delivery status received back through the live public tunnel --- end
to end, not simulated at any stage. The clinical-note extraction
pipeline was separately verified against real Epic sandbox patients'
\code{DocumentReference} notes, confirmed via a raw-data diagnostic
endpoint to correctly distinguish ``search failed'' from ``genuinely
zero notes'' for a given patient, and confirmed to correctly recognize
sandbox filler/test note content as content that legitimately yields no
structured extraction, rather than misreporting it as a pipeline defect.

\section{Discussion and limitations}

The principal limitation of this work is the same one the sister PFE
track states honestly for its own model: \textbf{the no-show model, and
every other ranking model in this pipeline, is trained on
research-guided synthetic data, not real clinic history}, because real
appointment and outreach outcome data was not available during this
internship. The synthetic dataset's probability structure follows
outpatient no-show literature rather than arbitrary labels, and the
model's feature-importance ranking (Table~\ref{tab:importance}) is
directionally consistent with that literature (recent booking behavior,
lead time, and specialty dominate, matching what published no-show
studies report as the strongest predictors) --- but real clinic
validation, with real subgroup behavior by provider, service, and
patient population, remains the necessary next step before any claim of
real-world predictive accuracy.

The causal holdout experiment's design is sound and enforced
structurally in the codebase, but its \emph{statistical power} --- how
quickly it can detect a real effect at a given confidence level ---
depends on real patient volume and real booking-rate variance that a
synthetic dataset cannot substitute for; this can only be assessed once
the system runs against a real clinic's ongoing patient population.

The clinical-note intelligence pipeline's practical accuracy is bounded
by the same limitation this project already documented honestly during
development: a meaningful share of the real Epic sandbox's own
clinical-note content is itself sandbox connectivity-test filler rather
than genuine documentation, which understates what extraction accuracy
against real clinical notes would actually look like.

Finally, several pieces of this platform's broader design are
explicitly out of scope as delivered work rather than claimed as done: a
production (non-sandbox) Epic connection for a real clinic,
commercial-scale WhatsApp Business template approval, and the
cross-module integration this platform's design anticipates with the
company's separate scheduling and billing/insurance products --- plus
the dental treatment-pathway prototype
(Section~\ref{sec:dental-prototype}), which remains deliberately
unreviewed and unshipped pending a named clinical advisor's sign-off.

\section{Conclusion}

This chapter presented the evaluation methodology, the no-show model's
genuinely recorded two-stage and feature-encoding-strategy results, the
guardrail and holdout-protection verification, the console migration's
end-to-end verification, and the project's honest limitations ---
chiefly, the reliance on research-guided synthetic training data pending
real clinic validation.

% =============================================================================
\chapter*{General conclusion}
\addcontentsline{toc}{chapter}{General conclusion}

This PFE internship focused on the design and implementation of Velo
Engage (VeloDoc), a rule-gated, machine-learning-ranked,
causally-measured patient-reactivation platform for outpatient clinics.
The project addressed a concrete operational gap distinct from the
sister PFE track's no-show/slot-recovery focus: not predicting whether
an \emph{already-booked} appointment will be attended, but deciding, from
a clinic's own dormant patient population, \emph{who should be
proactively contacted at all}, and proving that the outreach caused real
incremental bookings rather than merely reaching people who would have
returned regardless.

The proposed solution keeps Epic as the clinical source of truth
throughout, applies deterministic rule logic to every decision with a
real safety or compliance consequence (consent gating, contact
frequency, holdout-arm protection), and reserves machine learning for
what it is genuinely good at --- ranking within a space of options the
rule layer has already validated. Five trained models (no-show risk,
booking propensity, expected revenue, causal uplift, and predicted
time-to-book) combine into a single expected-value ranking, a
randomized holdout arm is enforced structurally rather than by
convention, and a separate clinical-note intelligence pipeline extracts
structured information from unstructured Epic notes under an explicit
content-safety gate.

The evaluation presented in Chapter~\ref{ch:evaluation} is honest about
what has and has not been verified: the no-show model's two-stage and
feature-encoding-strategy results are drawn directly from this
project's own recorded artifacts; the guardrail, holdout-protection, and
console-migration behavior were verified functionally and end to end
against the real running system, including a live Epic sandbox
connection and a live WhatsApp delivery round trip; and the remaining
four ranking models' final metrics, together with real-clinic validation
of the entire pipeline, remain the clearly identified next steps rather
than claimed results.

As future perspectives, the immediate priorities are: recalibrating
every ranking model against real clinic outcome data once a pilot
clinic's real Epic connection and real WhatsApp Business API deployment
exist; measuring the holdout experiment's actual statistical power
against real patient volume; and --- should the company's broader vision
proceed in this direction --- advancing the dental treatment-pathway
prototype only behind a named clinical advisor's review of its rule
graph, exactly as this report has scoped it throughout. A further
perspective, consistent with the company's own stated broader vision, is
integrating this platform with Velodoc's separate scheduling and
billing/insurance products, so a detected re-engagement opportunity,
once acted on, flows through to a real booked appointment and a real,
attributable revenue outcome rather than stopping at outreach alone.

% =============================================================================
\begin{thebibliography}{99}
\addcontentsline{toc}{chapter}{Bibliography}

% TODO: replace with your own real citations. The categories below match
% what this report actually draws on -- do not copy the sister report's
% list verbatim, since several of its sources are specific to
% no-show/waitlist modeling rather than to causal uplift, opportunity
% detection, or FHIR integration, which belong here instead.

\bibitem{noshow-lit} TODO: outpatient no-show prediction literature
(the same category the sister PFE report surveys).

\bibitem{catboost} Dorogush, A.\,V., Ershov, V., and Gulin, A. CatBoost:
gradient boosting with categorical features support. \emph{Workshop on
ML Systems at NIPS}, 2017.

\bibitem{shap} Lundberg, S.\,M. and Lee, S.-I. A unified approach to
interpreting model predictions. \emph{Advances in Neural Information
Processing Systems}, 2017.

\bibitem{uplift} TODO: T-learner / causal uplift modeling literature
(e.g., K\"unzel et al., ``Metalearners for estimating heterogeneous
treatment effects using machine learning'').

\bibitem{fhir} HL7 FHIR / SMART on FHIR specification. Available online:
\url{https://www.hl7.org/fhir/}.

\bibitem{temporal} Temporal documentation. Available online:
\url{https://docs.temporal.io}.

\bibitem{mlflow} MLflow documentation. Available online:
\url{https://mlflow.org}.

\bibitem{driftsurvey} Gama, J., Zliobaite, I., Bifet, A., Pechenizkiy,
M., and Bouchachia, A. A survey on concept drift adaptation. \emph{ACM
Computing Surveys}, 46(4), 44, 2014.

\bibitem{hiddendebt} Sculley, D., Holt, G., Golovin, D., et al. Hidden
technical debt in machine learning systems. \emph{Advances in Neural
Information Processing Systems}, 2015.

\end{thebibliography}

% =============================================================================
\appendix
\chapter{Appendices}

\section{Main API endpoint groups}

\begin{table}[h]
\centering
\caption{Main API endpoint groups}
\begin{tabular}{@{}p{5cm}p{8.5cm}@{}}
\toprule
Endpoint group & Purpose \\
\midrule
\code{/pull-cohort}, \code{/pull-bulk}, \code{/pull-bulk-sync}
(\code{ve_connect}) & Epic patient cohort ingestion (background and
synchronous variants). \\
\code{/patients/\{id\}/clinical-notes}, \code{/patients/\{id\}/raw}
(\code{ve_connect}) & Clinical-note retrieval and raw-FHIR diagnostic
access. \\
\code{/auth/login}, \code{/auth/callback} (\code{ve_connect}, proxied
via \code{ve_reach}) & Epic OAuth interactive flow. \\
\code{/api/v1/overview}, \code{/api/v1/opportunities},
\code{/api/v1/campaigns}, \code{/api/v1/campaigns/\{id\}}
(\code{ve_console}) & Console JSON API --- opportunities and campaign
detail. \\
\code{/api/v1/campaigns/\{id\}/outcome} (\code{ve_console}) & Manual
outcome marking (booked/attended), holdout-protected. \\
\code{/api/v1/whatsapp}, \code{/api/v1/whatsapp/messages/\{id\}}
(\code{ve_console}) & Conversation log. \\
\code{/api/v1/holdout}, \code{/api/v1/revenue} (\code{ve_console}) &
Causal experiment results and ROI (owner-gated). \\
\code{/api/v1/models}, \code{/api/v1/operations} (\code{ve_console}) &
Model registry health and pipeline operations (admin-gated). \\
\bottomrule
\end{tabular}
\end{table}

\section{Database separation}

\begin{table}[h]
\centering
\caption{Database roles}
\begin{tabular}{@{}p{3.5cm}p{10cm}@{}}
\toprule
Database concern & Role \\
\midrule
Clinic-scoped operational tables (Postgres) & Patients, features,
consent, opportunities, campaigns, outcomes, bookings, revenue ---
every table keyed by \code{clinic_id}. \\
Neo4j & The clinical-recall knowledge graph. \\
Redis & Ranking-time feature cache (short TTL, non-authoritative). \\
MLflow & Model artifacts, metrics, and registry pointers --- never
patient data. \\
\bottomrule
\end{tabular}
\end{table}

\section{Model artifacts}

\begin{table}[h]
\centering
\caption{Main model artifacts}
\begin{tabular}{@{}p{4.5cm}p{9cm}@{}}
\toprule
Artifact & Role \\
\midrule
No-show CatBoost/LightGBM model & Predicts appointment no-show risk,
feeding the expected-value ranking. \\
Booking propensity model & Predicts whether a candidate books at all if
contacted. \\
Expected-value (revenue) model & Predicts the monetary value of a
visit. \\
Uplift T-learner & Predicts the causal, incremental effect of contact
for a given patient. \\
Survival (time-to-book) model & Predicts time until booking if
contacted. \\
Lookalike/discovery model & Recommends a new opportunity family for a
patient based on similar patients. \\
WhatsApp intent classifier & Classifies inbound reply intent (booking,
confirmation, opt-out, ambiguous). \\
MedGemma + NemoGuard & Clinical-note structured extraction and its
content-safety gate. \\
\bottomrule
\end{tabular}
\end{table}

\section{Deployment commands}

\begin{verbatim}
docker compose -f infra/docker-compose.dev.yml up -d --build
docker compose -f infra/docker-compose.dev.yml logs -f ve_orchestrator
docker exec ve_temporal temporal schedule trigger \
  --address temporal:7233 --schedule-id daily-engagement-schedule-<clinic_id>
\end{verbatim}

\section{Code structure}

The main runtime modules are organized as one service per external
dependency or responsibility: \code{services/ve_connect} (Epic),
\code{services/ve_orchestrator} (Temporal pipeline),
\code{services/ve_reach} (WhatsApp/SMS), \code{services/ve_measure}
(booking/revenue attribution), \code{services/ve_console} (web
application), \code{services/ve_clinical_intel} (clinical-note
intelligence), \code{ml/} (training scripts and model registry), and
\code{policies/opportunities/} (the YAML opportunity-family rule
definitions).

\end{document}
